% preamble.tex — shared packages, macros, and theorem environments
% for the Quantum Scientific Computing lecture notes.

\usepackage[utf8]{inputenc}
\usepackage[T1]{fontenc}
\usepackage{lmodern}
\usepackage[margin=1in]{geometry}
\usepackage{amsmath,amssymb,amsthm,mathtools}
\usepackage{braket}          % Dirac notation: \ket, \bra, \braket, \bra{}\ket{}
\usepackage{bm}
\usepackage{graphicx}
\usepackage{booktabs}
\usepackage{enumitem}
\usepackage{xcolor}
\usepackage{tikz}
\usepackage[most]{tcolorbox}
\usepackage[colorlinks=true,linkcolor=blue,citecolor=blue,urlcolor=blue]{hyperref}
\usepackage[capitalize,noabbrev]{cleveref}

% ---------------------------------------------------------------------------
% Draft helpers (remove or disable before a "clean" build)
% ---------------------------------------------------------------------------
\newcommand{\todo}[1]{{\color{red}\textbf{[TODO: #1]}}}
\newcommand{\note}[1]{{\color{purple}\small(#1)}}

% ---------------------------------------------------------------------------
% Notation macros (kept consistent across all lectures)
% ---------------------------------------------------------------------------
\newcommand{\R}{\mathbb{R}}
\newcommand{\C}{\mathbb{C}}
\newcommand{\N}{\mathbb{N}}
\newcommand{\Z}{\mathbb{Z}}
\newcommand{\E}{\mathbb{E}}
\newcommand{\Prob}{\mathbb{P}}
\newcommand{\Hilb}{\mathcal{H}}
\newcommand{\ii}{\mathrm{i}}
\newcommand{\dd}{\mathrm{d}}
\newcommand{\ee}{\mathrm{e}}
\DeclareMathOperator{\tr}{tr}
\DeclareMathOperator{\Tr}{Tr}
\DeclareMathOperator{\Var}{Var}
\DeclareMathOperator{\Cov}{Cov}
\DeclareMathOperator{\rank}{rank}
\DeclareMathOperator{\diag}{diag}
\DeclareMathOperator{\poly}{poly}
\DeclareMathOperator{\polylog}{polylog}
\DeclareMathOperator*{\argmin}{arg\,min}
\DeclareMathOperator*{\argmax}{arg\,max}
\newcommand{\norm}[1]{\left\lVert #1 \right\rVert}
\newcommand{\abs}[1]{\left\lvert #1 \right\rvert}
\newcommand{\bigo}{\mathcal{O}}
\newcommand{\id}{\mathbb{I}}
\newcommand{\dg}{^{\dagger}}
\newcommand{\proj}[1]{\ket{#1}\!\bra{#1}}   % |x><x|

% Pauli / common gates
\newcommand{\Xgate}{\mathsf{X}}
\newcommand{\Ygate}{\mathsf{Y}}
\newcommand{\Zgate}{\mathsf{Z}}
\newcommand{\Hgate}{\mathsf{H}}

% ---------------------------------------------------------------------------
% Theorem environments
% ---------------------------------------------------------------------------
\theoremstyle{plain}
\newtheorem{theorem}{Theorem}[chapter]
\newtheorem{lemma}[theorem]{Lemma}
\newtheorem{proposition}[theorem]{Proposition}
\newtheorem{corollary}[theorem]{Corollary}
\theoremstyle{definition}
\newtheorem{definition}[theorem]{Definition}
\newtheorem{example}[theorem]{Example}
\newtheorem{exercise}{Exercise}[chapter]
\theoremstyle{remark}
\newtheorem{remark}[theorem]{Remark}

% Solution environment (unnumbered, italic "Solution." lead-in) — used only in solutions.tex.
\newenvironment{solution}{\begin{proof}[\normalfont\itshape Solution.]}{\end{proof}}

% ---------------------------------------------------------------------------
% Per-lecture structural helpers
% ---------------------------------------------------------------------------
% Learning-objectives box at the top of each lecture.
\newtcolorbox{objectives}{
  colback=blue!4, colframe=blue!45!black, title={Learning objectives},
  fonttitle=\bfseries, boxrule=0.5pt, arc=2pt, left=6pt, right=6pt, top=4pt, bottom=4pt}

% Uniform names for the closing sections of every lecture.
\newcommand{\exercisesection}{\section*{Exercises}\addcontentsline{toc}{section}{Exercises}}
\newcommand{\readingsection}{\section*{Further reading}\addcontentsline{toc}{section}{Further reading}}
